\section{Workflow: does it fit how I actually work?}
\label{sec:workflow}

The best safe tool fails if it does not fit. The eighth question asks whether the
system reduces burden or adds to it, and how it changes consent and the clinic day,
because a tool that fights the workflow will be worked around or abandoned.

[DRAFTING INSTRUCTIONS: in the full-framework, ground adoption in
\texttt{dietvorst2015algorithm} (algorithm aversion after seeing error) and
\texttt{novak2020patient} (patient stories and consent), and in the patient-journey
lineage \texttt{Kawchak2026CancerPatientsJourney}. Describe workflow integration,
burden measurement, and a verified consent process. Reproduce Mermaid
\texttt{adoption/mermaid/diagrams/11-clinician-workflow-journey.md},
\texttt{12-adoption-maturity-timeline.md}, and
\texttt{20-bedside-to-trial-trust-capstone.md} as colored cfwfig TikZ. Carry the
ADOPTION ACTIONS ``workflow integration'' and ``continuous monitoring''. Build a
body-width table mapping each of the eight questions to the mechanism that answers
it (the synthesis table).]

\begin{cfwfig}
\node[cfone] (a) at (0,0) {Bedside}; \node[cftwo] (b) at (3.0,0) {Site};
\node[cfact] (c) at (6.0,0) {Trial}; \node[cfhope] (d) at (9.0,0) {Platform trust};
\draw[cfedge] (a)--(b); \draw[cfedge] (b)--(c); \draw[cfedge] (c)--(d);
\end{cfwfig}
\vspace{-0.2cm}
\cfwcaption{Figure 9. From bedside to trial-wide trust (placeholder; expanded from
Mermaid 20 in the full-framework).}

Workflow is answered by fit: a system embedded in the clinic day and the consent
process, measured for burden, and scaled from one bedside decision to trial-wide
calibrated trust.
